\documentclass[11pt]{article}

\usepackage[margin=1in]{geometry}
\usepackage{amsmath,amssymb,amsthm,bm}
\usepackage{mathtools}
\usepackage{algorithm}
\usepackage{algpseudocode}
\usepackage{setspace}
\usepackage{graphicx}
\usepackage{booktabs}

\onehalfspacing

% Notation: bold lower-case = vectors, bold upper-case = matrices
\newcommand{\vect}[1]{\boldsymbol{#1}}
\newcommand{\mat}[1]{\mathbf{#1}}
\newcommand{\E}{\mathbb{E}}
\newcommand{\Var}{\mathbb{V}}
\newcommand{\R}{\mathbb{R}}
\newcommand{\IG}{\mathrm{Inv\text{-}Gamma}}
\newcommand{\N}{\mathcal{N}}

\title{Regularized Horseshoe Regression with Variational Bayes and Laplace--Delta Approximation}
\author{}
\date{}

\begin{document}
\maketitle

\section{Model Specification}

\subsection{Data and Likelihood}

Let
\begin{itemize}
  \item $n \in \mathbb{N}$ be the number of observations,
  \item $p \in \mathbb{N}$ be the number of predictors,
  \item $\vect{y} \in \R^n$ be the response vector,
  \item $\mat{X} \in \R^{n \times p}$ be the design matrix, with $i$-th row $\vect{x}_i^\top \in \R^{1 \times p}$,
  \item $\vect{\beta} \in \R^p$ be the regression coefficient vector,
  \item $\sigma^2 \in (0,\infty)$ be the observation noise variance.
\end{itemize}

The Gaussian linear regression model is
\begin{align}
  y_i \mid \vect{\beta}, \sigma^2 
  &\sim \N\big(\vect{x}_i^\top \vect{\beta}, \,\sigma^2\big),
  \quad i = 1,\dots,n, \\
  \vect{y} \mid \vect{\beta}, \sigma^2
  &\sim \N_n\big(\mat{X}\vect{\beta},\, \sigma^2 \mat{I}_n\big),
\end{align}
where $\mat{I}_n$ denotes the $n \times n$ identity matrix.

\subsection{Regularized Horseshoe Prior}

We introduce a global--local prior on $\vect{\beta}$ via the \emph{regularized horseshoe}. For $j=1,\dots,p$:

\begin{itemize}
  \item Local scale: $\lambda_j > 0$,
  \item Global scale: $\tau > 0$,
  \item Slab scale: $c^2 > 0$.
\end{itemize}

Define the prior variance of $\beta_j$ as
\begin{equation}
  V_j(\lambda_j, \tau, c^2)
  \;=\;
  \tau^2 \,\frac{c^2 \lambda_j^2}{c^2 + \tau^2 \lambda_j^2},
  \quad j = 1,\dots,p.
\end{equation}
Then
\begin{equation}
  \beta_j \mid \lambda_j,\tau,c^2 \;\sim\; \N\big(0,\, V_j(\lambda_j,\tau,c^2)\big),
  \quad j=1,\dots,p,
\end{equation}
independently across $j$.

The hierarchical prior on scales is:
\begin{align}
  \lambda_j &\sim C^+(0,1), 
  & j &= 1,\dots,p, \\
  \tau &\sim C^+(0,\tau_0^2), \\
  c^2 &\sim \IG\!\left(\frac{\nu}{2},\, \frac{\nu s^2}{2}\right),
\end{align}
where $C^+(0,a^2)$ denotes a half-Cauchy distribution with scale $a>0$, and $\IG(a,b)$ is the inverse-gamma distribution with density
\begin{equation}
  p(z) \propto z^{-a-1} \exp\!\left(-\frac{b}{z}\right), \quad z>0.
\end{equation}
The hyperparameters $\tau_0>0$, $\nu>0$, and $s>0$ encode prior information about sparsity and the slab scale.

For the noise variance we assume
\begin{equation}
  \sigma^2 \sim \IG(a_\sigma,b_\sigma), \qquad a_\sigma>0,\; b_\sigma>0.
\end{equation}

\subsection{Summary of Unknowns and Dimensions}

The full parameter vector is
\[
  \theta 
  = \big(\vect{\beta},\, \sigma^2,\, \lambda_1,\dots,\lambda_p,\, \tau,\, c^2\big),
\]
with dimensions:
\begin{itemize}
  \item $\vect{\beta} \in \R^p$,
  \item $\sigma^2 \in (0,\infty)$,
  \item $\lambda_j \in (0,\infty)$ for $j=1,\dots,p$,
  \item $\tau \in (0,\infty)$,
  \item $c^2 \in (0,\infty)$.
\end{itemize}

\section{Mean-Field Variational Approximation}

\subsection{Factorization}

We approximate the posterior $p(\theta \mid \vect{y})$ with a mean-field family
\begin{equation}
  q(\theta)
  \;=\;
  q_{\beta}(\vect{\beta})\, q_{\sigma}(\sigma^2)\, q_{\eta}(\vect{\eta}),
\end{equation}
where:
\begin{itemize}
  \item $q_\beta(\vect{\beta})$ approximates the marginal posterior of $\vect{\beta}$,
  \item $q_\sigma(\sigma^2)$ approximates the marginal posterior of $\sigma^2$,
  \item $q_\eta(\vect{\eta})$ approximates the joint posterior of the scale parameters,
\[
  \vect{\eta} \;=\; 
  \big(u_1,\dots,u_p,u_\tau,u_\kappa\big)^\top \in \R^{p+2}.
\]
\end{itemize}

We work in \emph{log-scale} for the non-conjugate block:
\begin{align}
  u_j &= \log \lambda_j,      & j &= 1,\dots,p, \\
  u_\tau &= \log \tau, \\
  u_\kappa &= \log c^2.
\end{align}
Thus,
\begin{align}
  \lambda_j &= e^{u_j},
  & \tau &= e^{u_\tau}, 
  & c^2 &= e^{u_\kappa}.
\end{align}

\subsection{Coordinate Ascent Variational Inference (CAVI)}

Let $\mathcal{L}(q)$ denote the evidence lower bound (ELBO)
\[
  \mathcal{L}(q) = \E_q[\log p(\vect{y},\theta)] - \E_q[\log q(\theta)].
\]
CAVI updates each factor by
\begin{align}
  \log q_\beta^*(\vect{\beta})
  &= \E_{q_\sigma q_\eta}[\log p(\vect{y},\theta)] + \text{const}, \\
  \log q_\sigma^*(\sigma^2)
  &= \E_{q_\beta q_\eta}[\log p(\vect{y},\theta)] + \text{const}, \\
  \log q_\eta^*(\vect{\eta})
  &= \E_{q_\beta q_\sigma}[\log p(\vect{y},\theta)] + \text{const}.
\end{align}

The first two updates are conjugate and yield closed forms. The third is non-conjugate; we approximate $q_\eta$ via a Laplace approximation in $\vect{\eta}$-space and use a Delta method for the required moments.

\section{Conjugate Variational Updates}

\subsection{Update for $q_\sigma(\sigma^2)$}

Given $q_\beta(\vect{\beta})$, the terms involving $\sigma^2$ are
\begin{align}
  \log p(\vect{y} \mid \vect{\beta},\sigma^2)
  &=
  -\frac{n}{2}\log(2\pi\sigma^2)
  -\frac{1}{2\sigma^2} \Vert \vect{y} - \mat{X}\vect{\beta} \Vert^2,\\
  \log p(\sigma^2)
  &=
  -(a_\sigma+1)\log \sigma^2 - \frac{b_\sigma}{\sigma^2} + \text{const}.
\end{align}
Taking expectation w.r.t.\ $q_\beta$ and collecting terms yields
\begin{equation}
  q_\sigma(\sigma^2)
  = \IG\big(\tilde{a}_\sigma, \tilde{b}_\sigma\big),
\end{equation}
with
\begin{align}
  \tilde{a}_\sigma 
  &= a_\sigma + \frac{n}{2},\\
  \tilde{b}_\sigma 
  &= b_\sigma + \frac{1}{2} S,
\end{align}
where
\begin{equation}
  S = \E_{q_\beta}\big[\Vert \vect{y} - \mat{X}\vect{\beta} \Vert^2\big]
    = \Vert \vect{y} - \mat{X}\vect{m}_\beta \Vert^2
      + \mathrm{tr}\big( \mat{X}\mat{\Sigma}_\beta \mat{X}^\top \big),
\end{equation}
and we denote
\[
  \vect{m}_\beta = \E_{q_\beta}[\vect{\beta}], 
  \quad \mat{\Sigma}_\beta = \Var_{q_\beta}(\vect{\beta}).
\]

We will need
\begin{equation}
  A := \E_{q_\sigma}\left[\frac{1}{\sigma^2}\right]
  = \frac{\tilde{a}_\sigma}{\tilde{b}_\sigma}.
\end{equation}

\subsection{Update for $q_\beta(\vect{\beta})$}

The terms involving $\vect{\beta}$ are
\begin{align}
  \log p(\vect{y} \mid \vect{\beta},\sigma^2)
  &= -\frac{n}{2}\log(2\pi\sigma^2)
     - \frac{1}{2\sigma^2}\Vert \vect{y} - \mat{X}\vect{\beta} \Vert^2,\\
  \log p(\vect{\beta} \mid \lambda_{1:p},\tau,c^2)
  &= -\frac12 \sum_{j=1}^p 
     \left( \log(2\pi V_j) + \frac{\beta_j^2}{V_j} \right),
\end{align}
where $V_j = V_j(\lambda_j,\tau,c^2)$.

Taking expectations w.r.t.\ $q_\sigma$ and $q_\eta$, define
\begin{equation}
  D_j := \E_{q_\eta} \left[\frac{1}{V_j(\lambda_j,\tau,c^2)}\right],
  \quad j=1,\dots,p.
\end{equation}
Then
\begin{equation}
  \log q_\beta^*(\vect{\beta})
  \propto
  -\frac{A}{2}\Vert \vect{y} - \mat{X}\vect{\beta} \Vert^2 
  -\frac12 \sum_{j=1}^p D_j \beta_j^2,
\end{equation}
which is the log-density of a multivariate Gaussian. Therefore,
\begin{equation}
  q_\beta(\vect{\beta})
  = \N_p\big(\vect{m}_\beta, \mat{\Sigma}_\beta\big),
\end{equation}
with
\begin{align}
  \mat{\Sigma}_\beta
  &= \left(
      A\,\mat{X}^\top \mat{X} 
      + \mathrm{diag}(D_1,\dots,D_p)
    \right)^{-1}, \\
  \vect{m}_\beta
  &= \mat{\Sigma}_\beta \, A\, \mat{X}^\top \vect{y}.
\end{align}

\section{Non-Conjugate Block and Laplace--Delta Approximation}

\subsection{Reparameterization and Basic Identities}

Recall:
\begin{align}
  \lambda_j &= e^{u_j},      & j &= 1,\dots,p, \\
  \tau      &= e^{u_\tau}, \\
  c^2       &= e^{u_\kappa}.
\end{align}
A key simplification arises from
\begin{equation}
  \frac{1}{V_j(\lambda_j,\tau,c^2)}
  = \frac{c^2 + \tau^2 \lambda_j^2}{\tau^2 c^2 \lambda_j^2}
  = \frac{1}{\tau^2 \lambda_j^2} + \frac{1}{c^2}.
\end{equation}
Define
\begin{align}
  T_j &= e^{-2u_j - 2u_\tau}, \\
  R   &= e^{-u_\kappa}, \\
  g_j &= T_j + R = \frac{1}{V_j(\boldsymbol{\eta})}.
\end{align}
Then $V_j(\boldsymbol{\eta}) = 1 / g_j$.

We also write
\begin{equation}
  S_j := \E_{q_\beta}[\beta_j^2]
       = m_{\beta,j}^2 + \Sigma_{\beta,jj},
  \quad j=1,\dots,p,
\end{equation}
where $m_{\beta,j}$ is the $j$-th component of $\vect{m}_\beta$ and $\Sigma_{\beta,jj}$ is the $j$-th diagonal entry of $\mat{\Sigma}_\beta$.

\subsection{Log-Target $f(\boldsymbol{\eta})$}

The optimal variational factor (up to a constant) is
\begin{equation}
  \log q_\eta^*(\vect{\eta})
  = \E_{q_\beta q_\sigma}[\log p(\vect{y},\theta)]
    + \log \left|\frac{\partial(\lambda_{1:p},\tau,c^2)}{\partial \vect{\eta}} \right|
    + \text{const}.
\end{equation}
After algebra and dropping constants independent of $\vect{\eta}$, we obtain
\begin{equation}
  f(\vect{\eta}) \;\propto\; \sum_{j=1}^p f_j(u_j,u_\tau,u_\kappa) 
  + f_{\mathrm{global}}(u_\tau,u_\kappa),
\end{equation}
with
\begin{align}
  f_j(u_j,u_\tau,u_\kappa)
  &=
  \frac12 \log g_j
  - \frac12 S_j g_j
  - \log\big(1+e^{2u_j}\big)
  + u_j, \label{eq:fj-def}\\[3pt]
  f_{\mathrm{global}}(u_\tau,u_\kappa)
  &=
  -\log\big(\tau_0^2 + e^{2u_\tau}\big)
  -\Big(\frac{\nu}{2}+1\Big)u_\kappa
  -\frac{\nu s^2}{2}e^{-u_\kappa}
  + u_\tau + u_\kappa. \label{eq:fglobal-def}
\end{align}
Here
\begin{equation}
  g_j 
  = e^{-2u_j-2u_\tau} + e^{-u_\kappa} 
  = T_j + R.
\end{equation}

Thus, explicitly,
\begin{align}
  f(\vect{\eta})
  &= \sum_{j=1}^p 
  \left[
    \frac12 \log\big(e^{-2u_j-2u_\tau} + e^{-u_\kappa}\big)
    - \frac12 S_j \big(e^{-2u_j-2u_\tau} + e^{-u_\kappa}\big)
    - \log\big(1+e^{2u_j}\big)
    + u_j
  \right] \nonumber\\
  &\quad -\log\big(\tau_0^2 + e^{2u_\tau}\big)
  -\Big(\frac{\nu}{2}+1\Big)u_\kappa
  -\frac{\nu s^2}{2}e^{-u_\kappa}
  + u_\tau + u_\kappa
  + \text{const}.
\end{align}

In the VB--LD scheme, we treat
\begin{equation}
  q_\eta(\vect{\eta}) \;\propto\; \exp\big(f(\vect{\eta})\big)
\end{equation}
as the target for the Laplace approximation.

\subsection{Gradient $\nabla f(\boldsymbol{\eta})$}

We now derive explicit expressions for $\nabla f(\vect{\eta})$.

\subsubsection*{Derivative with respect to $u_j$}

For a fixed $j$,
\[
  f_j(u_j,u_\tau,u_\kappa)
  = \frac12 \log g_j
  - \frac12 S_j g_j
  - \log(1+e^{2u_j})
  + u_j,
\]
with
\[
  g_j = T_j + R = e^{-2u_j-2u_\tau} + e^{-u_\kappa}.
\]
The derivative of $g_j$ w.r.t.\ $u_j$ is
\[
  \frac{\partial g_j}{\partial u_j}
  = -2 e^{-2u_j-2u_\tau}
  = -2 T_j.
\]
Then
\begin{align}
  \frac{\partial f_j}{\partial u_j}
  &= \frac12 \frac{1}{g_j}\frac{\partial g_j}{\partial u_j}
     - \frac12 S_j \frac{\partial g_j}{\partial u_j}
     - \frac{2 e^{2u_j}}{1+e^{2u_j}}
     + 1 \\
  &= \frac12 \frac{-2T_j}{g_j}
     - \frac12 S_j(-2T_j)
     - \frac{2 e^{2u_j}}{1+e^{2u_j}}
     + 1 \\
  &= -\frac{T_j}{g_j} + S_j T_j
     - \frac{2 e^{2u_j}}{1+e^{2u_j}} + 1.
\end{align}
No other term depends on $u_j$, so
\begin{equation}
  \boxed{
  \frac{\partial f}{\partial u_j}
  = -\frac{T_j}{g_j} + S_j T_j
    - \frac{2 e^{2u_j}}{1+e^{2u_j}} + 1,
  \quad j=1,\dots,p.
  }
\end{equation}

\subsubsection*{Derivative with respect to $u_\tau$}

For each $j$, $f_j$ depends on $u_\tau$ only through $g_j$:
\[
  \frac{\partial g_j}{\partial u_\tau}
  = -2 e^{-2u_j-2u_\tau}
  = -2T_j.
\]
Thus,
\begin{align}
  \frac{\partial f_j}{\partial u_\tau}
  &= \frac12 \frac{1}{g_j}\frac{\partial g_j}{\partial u_\tau}
     - \frac12 S_j \frac{\partial g_j}{\partial u_\tau}
  = -\frac{T_j}{g_j} + S_j T_j.
\end{align}

For the global term,
\[
  f_{\mathrm{global}}(u_\tau,u_\kappa)
  = -\log(\tau_0^2 + e^{2u_\tau})
    + u_\tau + \text{(terms independent of $u_\tau$)},
\]
so
\begin{align}
  \frac{\partial f_{\mathrm{global}}}{\partial u_\tau}
  &= -\frac{2 e^{2u_\tau}}{\tau_0^2 + e^{2u_\tau}} + 1.
\end{align}
Therefore
\begin{equation}
  \boxed{
  \frac{\partial f}{\partial u_\tau}
  = \sum_{j=1}^p \left(-\frac{T_j}{g_j} + S_j T_j\right)
    + 1 - \frac{2 e^{2u_\tau}}{\tau_0^2 + e^{2u_\tau}}.
  }
\end{equation}

\subsubsection*{Derivative with respect to $u_\kappa$}

Again
\[
  g_j = T_j + R, \quad R = e^{-u_\kappa},
\]
and
\[
  \frac{\partial g_j}{\partial u_\kappa}
  = -R.
\]
Thus, for each $j$,
\begin{align}
  \frac{\partial f_j}{\partial u_\kappa}
  &= \frac12 \frac{1}{g_j}\frac{\partial g_j}{\partial u_\kappa}
     - \frac12 S_j \frac{\partial g_j}{\partial u_\kappa} \\
  &= \frac12 \frac{-R}{g_j}
     - \frac12 S_j (-R) \\
  &= -\frac{R}{2g_j} + \frac{S_j R}{2}.
\end{align}

For the global term
\[
  f_{\mathrm{global}}(u_\tau,u_\kappa)
  = -\Big(\frac{\nu}{2}+1\Big)u_\kappa
    -\frac{\nu s^2}{2}e^{-u_\kappa}
    + u_\kappa + \text{(terms independent of $u_\kappa$)},
\]
we have
\begin{equation}
  \frac{\partial f_{\mathrm{global}}}{\partial u_\kappa}
  = -\Big(\frac{\nu}{2}+1\Big)
    + \frac{\nu s^2}{2}e^{-u_\kappa}
    + 1
  = -\frac{\nu}{2} + \frac{\nu s^2}{2}e^{-u_\kappa}.
\end{equation}

Summing over $j$:
\begin{equation}
  \boxed{
  \frac{\partial f}{\partial u_\kappa}
  = \sum_{j=1}^p\left(
      -\frac{R}{2g_j} + \frac{S_j R}{2}
    \right)
    -\frac{\nu}{2}
    + \frac{\nu s^2}{2}e^{-u_\kappa},
  }
\end{equation}
with $R = e^{-u_\kappa}$ and $g_j = e^{-2u_j-2u_\tau} + e^{-u_\kappa}$.

\subsection{Laplace Approximation for $q_\eta$}

Given the current $S_j$ (from $q_\beta$), the non-conjugate update is
\begin{equation}
  q_\eta(\vect{\eta})
  \propto \exp\big(f(\vect{\eta})\big).
\end{equation}

The Laplace approximation finds the mode
\begin{equation}
  \hat{\vect{\eta}} = \arg\max_{\vect{\eta} \in \R^{p+2}} f(\vect{\eta})
\end{equation}
using Newton iterations:
\begin{equation}
  \vect{\eta}^{(t+1)} 
  = \vect{\eta}^{(t)} - \left[\nabla^2 f(\vect{\eta}^{(t)})\right]^{-1} \nabla f(\vect{\eta}^{(t)}).
\end{equation}
At convergence, compute the Hessian $\nabla^2 f(\hat{\vect{\eta}})$ (e.g.\ via automatic differentiation) and set
\begin{equation}
  \mat{\Sigma}_\eta
  = -\left[\nabla^2 f(\hat{\vect{\eta}})\right]^{-1}.
\end{equation}
Then we approximate
\begin{equation}
  q_\eta(\vect{\eta})
  \approx \N\big(\hat{\vect{\eta}}, \mat{\Sigma}_\eta\big).
\end{equation}

\subsection{Delta Approximation for $D_j$}

The conjugate update for $q_\beta$ requires
\begin{equation}
  D_j = \E_{q_\eta}\left[\frac{1}{V_j(\vect{\eta})}\right]
      = \E_{q_\eta}[g_j(\vect{\eta})],
\end{equation}
where
\begin{equation}
  g_j(\vect{\eta}) = \frac{1}{V_j(\vect{\eta})}
  = e^{-2u_j-2u_\tau} + e^{-u_\kappa} = T_j + R,
\end{equation}
with
\begin{equation}
  T_j = e^{-2u_j-2u_\tau}, 
  \qquad
  R   = e^{-u_\kappa}.
\end{equation}

Under the Gaussian approximation $q_\eta \approx \N(\hat{\vect{\eta}},\mat{\Sigma}_\eta)$, the second-order Delta method gives
\begin{equation}
  D_j \;\approx\;
  g_j(\hat{\vect{\eta}})
  + \frac12 \mathrm{tr}\big( \nabla^2 g_j(\hat{\vect{\eta}})\, \mat{\Sigma}_\eta \big),
\end{equation}
where $\nabla^2 g_j$ denotes the Hessian of $g_j$ with respect to the log-scale parameters $\vect{\eta}$.

The function $g_j$ depends only on the three components $(u_j,u_\tau,u_\kappa)$:
\[
  g_j(u_j,u_\tau,u_\kappa) = e^{-2u_j-2u_\tau} + e^{-u_\kappa},
\]
so the Hessian $\nabla^2 g_j$ is zero outside the $3\times3$ block corresponding to $(u_j,u_\tau,u_\kappa)$. Ordering these coordinates as $(u_j,u_\tau,u_\kappa)$, the non-zero block is
\begin{equation}
  \nabla^2 g_j(u_j,u_\tau,u_\kappa)
  =
  \begin{pmatrix}
    \displaystyle \frac{\partial^2 g_j}{\partial u_j^2} &
    \displaystyle \frac{\partial^2 g_j}{\partial u_j \partial u_\tau} &
    \displaystyle \frac{\partial^2 g_j}{\partial u_j \partial u_\kappa} \\[6pt]
    \displaystyle \frac{\partial^2 g_j}{\partial u_\tau \partial u_j} &
    \displaystyle \frac{\partial^2 g_j}{\partial u_\tau^2} &
    \displaystyle \frac{\partial^2 g_j}{\partial u_\tau \partial u_\kappa} \\[6pt]
    \displaystyle \frac{\partial^2 g_j}{\partial u_\kappa \partial u_j} &
    \displaystyle \frac{\partial^2 g_j}{\partial u_\kappa \partial u_\tau} &
    \displaystyle \frac{\partial^2 g_j}{\partial u_\kappa^2}
  \end{pmatrix}
  =
  \begin{pmatrix}
    4T_j & 4T_j & 0 \\
    4T_j & 4T_j & 0 \\
    0    & 0    & R
  \end{pmatrix},
\end{equation}
i.e.
\begin{equation}
  \nabla^2 g_j(u_j,u_\tau,u_\kappa)
  =
  \begin{pmatrix}
    4 e^{-2u_j-2u_\tau} & 4 e^{-2u_j-2u_\tau} & 0 \\
    4 e^{-2u_j-2u_\tau} & 4 e^{-2u_j-2u_\tau} & 0 \\
    0                   & 0                   & e^{-u_\kappa}
  \end{pmatrix}.
\end{equation}
Therefore, in the Delta approximation, only this $3\times3$ block of $\nabla^2 g_j(\hat{\vect{\eta}})$ and the corresponding $3\times3$ block of the covariance matrix $\mat{\Sigma}_\eta$ contribute to the trace term in
\[
  D_j \;\approx\;
  g_j(\hat{\vect{\eta}})
  + \frac12 \mathrm{tr}\big( \nabla^2 g_j(\hat{\vect{\eta}})\, \mat{\Sigma}_\eta \big).
\]

\section{VB--LD Algorithm (Compact Description)}

\begin{algorithm}[h]
\caption{VB with Laplace--Delta (VB--LD) for Regularized Horseshoe Regression}
\begin{algorithmic}[1]
\State Initialize $q_\beta^{(0)}(\vect{\beta}) = \N(\vect{m}_\beta^{(0)},\mat{\Sigma}_\beta^{(0)})$, 
       $q_\sigma^{(0)}(\sigma^2) = \IG(\tilde{a}_\sigma^{(0)},\tilde{b}_\sigma^{(0)})$,
       and $\hat{\vect{\eta}}^{(0)}$.
\For{$k = 0,1,2,\dots$ until convergence}
  \State \textbf{(Delta step)} Given $q_\eta^{(k)} \approx \N(\hat{\vect{\eta}}^{(k)},\mat{\Sigma}_\eta^{(k)})$, compute
    \[
      D_j^{(k)} \approx g_j(\hat{\vect{\eta}}^{(k)}) 
      + \frac12 \mathrm{tr}\big(\nabla^2 g_j(\hat{\vect{\eta}}^{(k)}) \mat{\Sigma}_\eta^{(k)}\big),
      \quad j=1,\dots,p.
    \]
  \State \textbf{(Update $q_\beta$)} Compute
    \[
      A^{(k)} = \frac{\tilde{a}_\sigma^{(k)}}{\tilde{b}_\sigma^{(k)}},
    \]
    and update
    \[
      \mat{\Sigma}_\beta^{(k+1)} 
      = \left(A^{(k)}\mat{X}^\top\mat{X} + \mathrm{diag}(D_1^{(k)},\dots,D_p^{(k)})\right)^{-1},
    \]
    \[
      \vect{m}_\beta^{(k+1)} 
      = \mat{\Sigma}_\beta^{(k+1)} A^{(k)} \mat{X}^\top \vect{y}.
    \]
  \State \textbf{(Update $q_\sigma$)} Compute
    \[
      S^{(k+1)} = \Vert \vect{y} - \mat{X}\vect{m}_\beta^{(k+1)} \Vert^2
      + \mathrm{tr}\big(\mat{X}\mat{\Sigma}_\beta^{(k+1)}\mat{X}^\top\big),
    \]
    and set
    \[
      \tilde{a}_\sigma^{(k+1)} = a_\sigma + \frac{n}{2}, \quad
      \tilde{b}_\sigma^{(k+1)} = b_\sigma + \frac{S^{(k+1)}}{2}.
    \]
  \State \textbf{(Laplace step for $q_\eta$)} 
    \Statex \quad (a) Compute $S_j^{(k+1)} = m_{\beta,j}^{(k+1)2} + \Sigma_{\beta,jj}^{(k+1)}$ for $j=1,\dots,p$.
    \Statex \quad (b) Define $f(\vect{\eta})$ using $S_j^{(k+1)}$ and gradients $\nabla f(\vect{\eta})$ as above.
    \Statex \quad (c) Use Newton iterations to find 
      $\displaystyle \hat{\vect{\eta}}^{(k+1)} = \arg\max_{\vect{\eta}} f(\vect{\eta})$.
    \Statex \quad (d) Compute 
      $\mat{\Sigma}_\eta^{(k+1)} = -\big[\nabla^2 f(\hat{\vect{\eta}}^{(k+1)})\big]^{-1}$.
    \Statex \quad (e) Set $q_\eta^{(k+1)}(\vect{\eta}) \approx \N(\hat{\vect{\eta}}^{(k+1)},\mat{\Sigma}_\eta^{(k+1)})$.
\EndFor
\end{algorithmic}
\end{algorithm}

This algorithm alternates between:
\begin{itemize}
  \item Conjugate updates for $(q_\beta,q_\sigma)$, given approximate moments $D_j$,
  \item A Laplace approximation for the non-conjugate scale block $q_\eta$,
  \item A Delta-method correction for the expectations $D_j = \E_{q_\eta}[1/V_j]$.
\end{itemize}

\section{Evidence Lower Bound (ELBO)}

For the mean-field variational family
\[
  q(\theta) = q_\beta(\vect{\beta})\,q_\sigma(\sigma^2)\,q_\eta(\vect{\eta}),
\]
the evidence lower bound (ELBO) is
\begin{equation}
  \mathcal{L}(q)
  = \E_q[\log p(\vect{y},\theta)] - \E_q[\log q(\theta)].
\end{equation}
Using the model and factorization above, and working up to an additive constant independent of the variational parameters, we can write
\begin{equation}
  \mathcal{L}(q)
  \approx
  \underbrace{\E_{q_\beta,q_\sigma}[\log p(\vect{y}\mid\vect{\beta},\sigma^2)]}_{\text{likelihood}}
  + \underbrace{\E_{q_\sigma}[\log p(\sigma^2)]}_{\text{noise prior}}
  + \underbrace{\E_{q_\eta}[f(\vect{\eta})]}_{\text{scale block}}
  - \underbrace{\E_{q_\beta}[\log q_\beta]}_{\text{entropy of }q_\beta}
  - \underbrace{\E_{q_\sigma}[\log q_\sigma]}_{\text{entropy of }q_\sigma}
  - \underbrace{\E_{q_\eta}[\log q_\eta]}_{\text{entropy of }q_\eta},
\end{equation}
where $f(\vect{\eta})$ is the non-conjugate log-target defined in the previous section.

\subsection*{Components in Closed Form}

Recall
\[
  q_\beta(\vect{\beta}) = \N_p(\vect{m}_\beta,\mat{\Sigma}_\beta),
  \qquad
  q_\sigma(\sigma^2) = \IG(\tilde a_\sigma,\tilde b_\sigma),
  \qquad
  q_\eta(\vect{\eta}) \approx \N_{p+2}(\hat{\vect{\eta}},\mat{\Sigma}_\eta),
\]
with
\begin{equation}
  A := \E_{q_\sigma}\!\left[\frac{1}{\sigma^2}\right]
  = \frac{\tilde a_\sigma}{\tilde b_\sigma},
  \qquad
  S := \E_{q_\beta}\big[\Vert \vect{y} - \mat{X}\vect{\beta}\Vert^2\big]
  = \Vert \vect{y} - \mat{X}\vect{m}_\beta\Vert^2
  + \mathrm{tr}\big(\mat{X}\mat{\Sigma}_\beta\mat{X}^\top\big).
\end{equation}
For the inverse-gamma,
\begin{equation}
  \E_{q_\sigma}[\log \sigma^2] = \log \tilde b_\sigma - \psi(\tilde a_\sigma),
\end{equation}
where $\psi(\cdot)$ is the digamma function.

\paragraph{Likelihood term.}
\[
  \log p(\vect{y}\mid\vect{\beta},\sigma^2)
  = -\frac{n}{2}\log(2\pi)
    -\frac{n}{2}\log\sigma^2
    -\frac{1}{2\sigma^2}\Vert \vect{y}-\mat{X}\vect{\beta}\Vert^2,
\]
so
\begin{equation}
  \E_{q_\beta,q_\sigma}[\log p(\vect{y}\mid\vect{\beta},\sigma^2)]
  = -\frac{n}{2}\log(2\pi)
    -\frac{n}{2}\big(\log \tilde b_\sigma - \psi(\tilde a_\sigma)\big)
    -\frac{1}{2} A S.
\end{equation}

\paragraph{Noise prior term.}
For $\sigma^2 \sim \IG(a_\sigma,b_\sigma)$ with density
\(
  p(\sigma^2) \propto (\sigma^2)^{-a_\sigma-1}\exp(-b_\sigma/\sigma^2),
\)
\[
  \log p(\sigma^2)
  = a_\sigma \log b_\sigma - \log\Gamma(a_\sigma)
    - (a_\sigma+1)\log\sigma^2
    - \frac{b_\sigma}{\sigma^2},
\]
hence
\begin{equation}
  \E_{q_\sigma}[\log p(\sigma^2)]
  = a_\sigma \log b_\sigma - \log\Gamma(a_\sigma)
    - (a_\sigma+1)\big(\log \tilde b_\sigma - \psi(\tilde a_\sigma)\big)
    - b_\sigma A.
\end{equation}

\paragraph{Scale-block term.}
With
\[
  f(\vect{\eta})
  = \sum_{j=1}^p f_j(u_j,u_\tau,u_\kappa)
    + f_{\mathrm{global}}(u_\tau,u_\kappa),
\]
defined earlier, we approximate $\E_{q_\eta}[f(\vect{\eta})]$ by a Laplace--Delta expansion around the mode $\hat{\vect{\eta}}$:
\begin{equation}
  \E_{q_\eta}[f(\vect{\eta})]
  \;\approx\;
  f(\hat{\vect{\eta}})
  + \frac12\,\mathrm{tr}\!\big(\nabla^2 f(\hat{\vect{\eta}})\,\mat{\Sigma}_\eta\big),
\end{equation}
where $\nabla^2 f(\hat{\vect{\eta}})$ is the Hessian used in the Laplace approximation and $\mat{\Sigma}_\eta = -[\nabla^2 f(\hat{\vect{\eta}})]^{-1}$.

\paragraph{Entropy terms.}

For the Gaussian factor $q_\beta$,
\begin{equation}
  \E_{q_\beta}[\log q_\beta]
  = -\frac12\big(p\log(2\pi) + \log|\mat{\Sigma}_\beta| + p\big),
  \quad\Rightarrow\quad
  -\E_{q_\beta}[\log q_\beta]
  = \frac12\big(p\log(2\pi e) + \log|\mat{\Sigma}_\beta|\big).
\end{equation}

For the inverse-gamma factor $q_\sigma(\sigma^2) = \IG(\tilde a_\sigma,\tilde b_\sigma)$,
\begin{align}
  \E_{q_\sigma}[\log q_\sigma]
  &= \tilde a_\sigma \log \tilde b_\sigma - \log\Gamma(\tilde a_\sigma)
    - (\tilde a_\sigma+1)\big(\log \tilde b_\sigma - \psi(\tilde a_\sigma)\big)
    - \tilde b_\sigma A, \\
  -\E_{q_\sigma}[\log q_\sigma]
  &= -\tilde a_\sigma \log \tilde b_\sigma + \log\Gamma(\tilde a_\sigma)
    + (\tilde a_\sigma+1)\big(\log \tilde b_\sigma - \psi(\tilde a_\sigma)\big)
    + \tilde b_\sigma A.
\end{align}

For the Gaussian factor $q_\eta(\vect{\eta}) \approx \N_{p+2}(\hat{\vect{\eta}},\mat{\Sigma}_\eta)$,
\begin{equation}
  \E_{q_\eta}[\log q_\eta]
  = -\frac12\big((p+2)\log(2\pi) + \log|\mat{\Sigma}_\eta| + (p+2)\big),
  \quad\Rightarrow\quad
  -\E_{q_\eta}[\log q_\eta]
  = \frac12\big((p+2)\log(2\pi e) + \log|\mat{\Sigma}_\eta|\big).
\end{equation}

\subsection*{Practical ELBO for Monitoring Convergence}

Collecting the pieces and omitting additive constants independent of
$(\vect{m}_\beta,\mat{\Sigma}_\beta,\tilde a_\sigma,\tilde b_\sigma,\hat{\vect{\eta}},\mat{\Sigma}_\eta)$, an implementable approximation to the ELBO is
\begin{align}
  \mathcal{L}(q)
  \approx\;
  &-\frac{n}{2}\big(\log \tilde b_\sigma - \psi(\tilde a_\sigma)\big)
   -\frac{1}{2} A S
   -(a_\sigma+1)\big(\log \tilde b_\sigma - \psi(\tilde a_\sigma)\big)
   - b_\sigma A \nonumber\\
  &+ f(\hat{\vect{\eta}})
   + \frac12\,\mathrm{tr}\!\big(\nabla^2 f(\hat{\vect{\eta}})\,\mat{\Sigma}_\eta\big) \nonumber\\
  &+ \frac12\big(p\log(2\pi e) + \log|\mat{\Sigma}_\beta|\big)
   -\E_{q_\sigma}[\log q_\sigma]
   + \frac12\big((p+2)\log(2\pi e) + \log|\mat{\Sigma}_\eta|\big)
   \;+\; \text{const},
\end{align}
where $A$, $S$, $f(\cdot)$, $\hat{\vect{\eta}}$, and $\mat{\Sigma}_\eta$ are as defined above. In practice, one may drop the explicit constants and track only relative changes in $\mathcal{L}(q)$ as a stopping criterion for the VB--LD iterations.

\section{Regularized Horseshoe Scale Block under exAL and AL Likelihoods}

\subsection{Setup and notational conventions}

This section records the full joint expressions for the latent regularized horseshoe (RHS)
scale parameters under the extended asymmetric Laplace (exAL) likelihood and its
asymmetric Laplace (AL) special case.

To avoid clashes with notation already used in the document, we adopt the following
conventions only within this section:
\begin{itemize}
  \item the quantile level is denoted by $q \in (0,1)$ (rather than $p$, which is already used above for the number of predictors),
  \item the slab hyperparameter in the prior for $c^2$ is denoted by $s_0>0$ (rather than $s$, which is reserved below for the exAL latent variables),
  \item the exAL/AL hierarchy is written in terms of a positive scale parameter $\sigma>0$.
\end{itemize}

Let $(y_i,\vect{x}_i)$, $i=1,\dots,n$, denote the response--covariate pairs, with
$\vect{x}_i \in \R^p$, and let
\begin{equation}
  Q_q(y_i \mid \vect{x}_i) = \vect{x}_i^\top \vect{\beta},
  \qquad i=1,\dots,n,
\end{equation}
so that $\vect{x}_i^\top \vect{\beta}$ is the conditional $q$-th quantile of $y_i$.

The RHS prior on the regression coefficients is
\begin{equation}
  \beta_j \mid \lambda_j,\tau,c^2
  \sim \N\!\big(0, V_j(\lambda_j,\tau,c^2)\big),
  \qquad j=1,\dots,p,
\end{equation}
with
\begin{equation}
  V_j(\lambda_j,\tau,c^2)
  =
  \tau^2 \,\frac{c^2 \lambda_j^2}{c^2 + \tau^2 \lambda_j^2},
  \qquad j=1,\dots,p.
\end{equation}
The scale hyperpriors are
\begin{align}
  \lambda_j &\sim C^+(0,1), \qquad j=1,\dots,p,\\
  \tau &\sim C^+(0,\tau_0),\\
  c^2 &\sim \IG\!\left(\frac{\nu}{2}, \frac{\nu s_0^2}{2}\right),
\end{align}
so that, up to normalizing constants,
\begin{align}
  p(\lambda_j) &\propto \frac{1}{1+\lambda_j^2}, \qquad \lambda_j>0,\\
  p(\tau) &\propto \frac{1}{1+\tau^2/\tau_0^2}, \qquad \tau>0,\\
  p(c^2) &\propto (c^2)^{-(\nu/2+1)}
  \exp\!\left(-\frac{\nu s_0^2}{2c^2}\right), \qquad c^2>0.
\end{align}

If an intercept is left unshrunk in the final model, then all products and sums below
should simply be taken over the coefficients that are assigned the RHS prior.

\subsection{Full joint density under the exAL likelihood}

For the exAL likelihood, introduce the per-observation latent variables
\[
  v_i > 0, \qquad s_i > 0, \qquad i=1,\dots,n,
\]
and define
\begin{align}
  m_i(\vect{\beta},\sigma,\gamma,v_i,s_i)
  &=
  \vect{x}_i^\top \vect{\beta}
  + \sigma C(\gamma)\lvert \gamma \rvert s_i
  + A(\gamma) v_i,\\
  \omega_i^2(\sigma,\gamma,v_i)
  &=
  \sigma B(\gamma) v_i.
\end{align}
Then the exAL hierarchical representation is
\begin{align}
  y_i \mid \vect{\beta},\sigma,\gamma,v_i,s_i
  &\sim \N\!\big(m_i(\vect{\beta},\sigma,\gamma,v_i,s_i),\, \omega_i^2(\sigma,\gamma,v_i)\big),\\
  v_i \mid \sigma
  &\sim \mathrm{Exp}(\text{mean}=\sigma),\\
  s_i
  &\sim \N^+(0,1),
\end{align}
independently across $i=1,\dots,n$.

Writing
\[
  \vect{v}=(v_1,\dots,v_n)^\top,
  \qquad
  \vect{s}=(s_1,\dots,s_n)^\top,
  \qquad
  \vect{\lambda}=(\lambda_1,\dots,\lambda_p)^\top,
\]
the full joint density under the exAL likelihood factorizes as
\begin{align}
  &p(\vect{y},\vect{v},\vect{s},\vect{\beta},\vect{\lambda},\tau,c^2,\sigma,\gamma)
  \nonumber\\
  &\qquad=
  \left\{
    \prod_{i=1}^n
    p(y_i \mid \vect{\beta},\sigma,\gamma,v_i,s_i)
  \right\}
  \left\{
    \prod_{i=1}^n
    p(v_i \mid \sigma)\,p(s_i)
  \right\}
  p(\vect{\beta}\mid \vect{\lambda},\tau,c^2)\,
  p(\vect{\lambda})\,p(\tau)\,p(c^2)\,p(\sigma)\,p(\gamma).
\end{align}

More explicitly,
\begin{align}
  p(y_i \mid \vect{\beta},\sigma,\gamma,v_i,s_i)
  &=
  \frac{1}{\sqrt{2\pi\,\sigma B(\gamma) v_i}}
  \exp\!\left[
    -\frac{
      \big\{
        y_i - \vect{x}_i^\top \vect{\beta}
        - \sigma C(\gamma)\lvert\gamma\rvert s_i
        - A(\gamma) v_i
      \big\}^2
    }{
      2\sigma B(\gamma) v_i
    }
  \right],\\
  p(v_i\mid \sigma)
  &=
  \frac{1}{\sigma}\exp\!\left(-\frac{v_i}{\sigma}\right),
  \qquad v_i>0,\\
  p(s_i)
  &=
  2\phi(s_i),
  \qquad s_i>0,\\
  p(\vect{\beta}\mid \vect{\lambda},\tau,c^2)
  &=
  \prod_{j=1}^p
  \frac{1}{\sqrt{2\pi V_j(\lambda_j,\tau,c^2)}}
  \exp\!\left(
    -\frac{\beta_j^2}{2V_j(\lambda_j,\tau,c^2)}
  \right).
\end{align}

Therefore, up to an additive constant independent of
$(\vect{\beta},\vect{\lambda},\tau,c^2,\sigma,\gamma,\vect{v},\vect{s})$, the exAL log-joint density is
\begin{align}
  &\log p(\vect{y},\vect{v},\vect{s},\vect{\beta},\vect{\lambda},\tau,c^2,\sigma,\gamma)
  \nonumber\\
  &\qquad=
  -\frac{1}{2}\sum_{i=1}^n
  \left[
    \log\!\big(\sigma B(\gamma) v_i\big)
    +
    \frac{
      \big\{
        y_i - \vect{x}_i^\top \vect{\beta}
        - \sigma C(\gamma)\lvert\gamma\rvert s_i
        - A(\gamma) v_i
      \big\}^2
    }{
      \sigma B(\gamma) v_i
    }
  \right]
  \nonumber\\
  &\qquad\quad
  - \sum_{i=1}^n \left(\log \sigma + \frac{v_i}{\sigma}\right)
  - \frac{1}{2}\sum_{i=1}^n s_i^2
  \nonumber\\
  &\qquad\quad
  - \frac{1}{2}\sum_{j=1}^p
  \left[
    \log V_j(\lambda_j,\tau,c^2)
    +
    \frac{\beta_j^2}{V_j(\lambda_j,\tau,c^2)}
  \right]
  \nonumber\\
  &\qquad\quad
  - \sum_{j=1}^p \log(1+\lambda_j^2)
  - \log\!\left(1+\frac{\tau^2}{\tau_0^2}\right)
  - \left(\frac{\nu}{2}+1\right)\log c^2
  - \frac{\nu s_0^2}{2c^2}
  \nonumber\\
  &\qquad\quad
  + \log p(\sigma) + \log p(\gamma).
  \label{eq:logjoint-exAL-RHS}
\end{align}

\subsection{Exact conditional kernel for the RHS scale block under exAL}

The latent RHS scale block is $(\vect{\lambda},\tau,c^2)$. Conditional on $\vect{\beta}$,
all exAL likelihood terms and all mixture-latent terms are constant with respect to
$(\vect{\lambda},\tau,c^2)$. Hence,
\begin{equation}
  p(\vect{\lambda},\tau,c^2 \mid \vect{\beta},\vect{y},\vect{v},\vect{s},\sigma,\gamma)
  \propto
  p(\vect{\beta}\mid \vect{\lambda},\tau,c^2)\,
  p(\vect{\lambda})\,p(\tau)\,p(c^2).
\end{equation}
Therefore,
\begin{align}
  &p(\vect{\lambda},\tau,c^2 \mid \vect{\beta},\vect{y},\vect{v},\vect{s},\sigma,\gamma)
  \nonumber\\
  &\qquad\propto
  \left\{
    \prod_{j=1}^p
    V_j(\lambda_j,\tau,c^2)^{-1/2}
    \exp\!\left(
      -\frac{\beta_j^2}{2V_j(\lambda_j,\tau,c^2)}
    \right)
    (1+\lambda_j^2)^{-1}
  \right\}
  \left(1+\frac{\tau^2}{\tau_0^2}\right)^{-1}
  \nonumber\\
  &\qquad\quad\times
  (c^2)^{-(\nu/2+1)}
  \exp\!\left(-\frac{\nu s_0^2}{2c^2}\right).
  \label{eq:rhs-kernel-exAL}
\end{align}
Equivalently, up to an additive constant,
\begin{align}
  &\log p(\vect{\lambda},\tau,c^2 \mid \vect{\beta},\vect{y},\vect{v},\vect{s},\sigma,\gamma)
  \nonumber\\
  &\qquad=
  - \frac{1}{2}\sum_{j=1}^p
  \left[
    \log V_j(\lambda_j,\tau,c^2)
    + \frac{\beta_j^2}{V_j(\lambda_j,\tau,c^2)}
  \right]
  - \sum_{j=1}^p \log(1+\lambda_j^2)
  \nonumber\\
  &\qquad\quad
  - \log\!\left(1+\frac{\tau^2}{\tau_0^2}\right)
  - \left(\frac{\nu}{2}+1\right)\log c^2
  - \frac{\nu s_0^2}{2c^2}.
  \label{eq:rhs-logkernel-exAL}
\end{align}

A useful identity is
\begin{equation}
  \frac{1}{V_j(\lambda_j,\tau,c^2)}
  =
  \frac{c^2+\tau^2\lambda_j^2}{\tau^2 c^2 \lambda_j^2}
  =
  \frac{1}{\tau^2\lambda_j^2} + \frac{1}{c^2}.
  \label{eq:Vinv-identity}
\end{equation}

\subsection{Log-scale representation of the RHS block}

For Laplace--Delta calculations it is convenient to work on the unconstrained coordinates
\begin{align}
  u_j &= \log \lambda_j, \qquad j=1,\dots,p,\\
  u_\tau &= \log \tau,\\
  u_\kappa &= \log c^2,
\end{align}
and to collect them as
\[
  \vect{\eta}
  =
  (u_1,\dots,u_p,u_\tau,u_\kappa)^\top \in \R^{p+2}.
\]
Under this transformation,
\begin{equation}
  \lambda_j = e^{u_j},
  \qquad
  \tau = e^{u_\tau},
  \qquad
  c^2 = e^{u_\kappa},
\end{equation}
and the Jacobian determinant is
\begin{equation}
  \left|
    \frac{\partial(\lambda_1,\dots,\lambda_p,\tau,c^2)}{\partial \vect{\eta}}
  \right|
  =
  \exp\!\left(\sum_{j=1}^p u_j + u_\tau + u_\kappa\right).
\end{equation}

Define
\begin{align}
  g_j(\vect{\eta})
  &:=
  e^{-2u_j-2u_\tau} + e^{-u_\kappa},
  \qquad j=1,\dots,p.
\end{align}
By \eqref{eq:Vinv-identity},
\begin{equation}
  g_j(\vect{\eta})
  =
  \frac{1}{V_j(\lambda_j,\tau,c^2)},
  \qquad
  V_j(\lambda_j,\tau,c^2)=\frac{1}{g_j(\vect{\eta})}.
\end{equation}

Hence the exact transformed log-kernel for the RHS block is
\begin{align}
  f_{\mathrm{RHS}}(\vect{\eta}\mid \vect{\beta})
  &=
  \log p(\vect{\lambda},\tau,c^2\mid \vect{\beta})
  +
  \log\left|
    \frac{\partial(\lambda_1,\dots,\lambda_p,\tau,c^2)}{\partial \vect{\eta}}
  \right|
  + \text{const}
  \nonumber\\
  &=
  \sum_{j=1}^p
  \left[
    \frac{1}{2}\log g_j(\vect{\eta})
    - \frac{1}{2}\beta_j^2 g_j(\vect{\eta})
    - \log\!\big(1+e^{2u_j}\big)
    + u_j
  \right]
  \nonumber\\
  &\quad
  - \log\!\big(\tau_0^2 + e^{2u_\tau}\big)
  - \frac{\nu}{2}u_\kappa
  - \frac{\nu s_0^2}{2}e^{-u_\kappa}
  + u_\tau
  + \text{const}.
  \label{eq:rhs-logtarget-transformed}
\end{align}

For variational Bayes, if one needs the corresponding expected target under a Gaussian
factor $q_\beta(\vect{\beta})$, one simply replaces $\beta_j^2$ in
\eqref{eq:rhs-logtarget-transformed} by
\begin{equation}
  S_j := \E_{q_\beta}[\beta_j^2],
  \qquad j=1,\dots,p.
\end{equation}

\subsection{AL likelihood as a special case}

At a fixed quantile level $q \in (0,1)$, define the standard AL constants
\begin{equation}
  A_q = \frac{1-2q}{q(1-q)},
  \qquad
  B_q = \frac{2}{q(1-q)}.
\end{equation}
Under the AL working likelihood,
\begin{align}
  y_i \mid \vect{\beta},\sigma,v_i
  &\sim
  \N\!\big(\vect{x}_i^\top \vect{\beta} + A_q v_i,\,
  \sigma B_q v_i\big),\\
  v_i \mid \sigma
  &\sim \mathrm{Exp}(\text{mean}=\sigma),
\end{align}
independently across $i=1,\dots,n$.

The full joint density under the AL likelihood is then
\begin{align}
  &p(\vect{y},\vect{v},\vect{\beta},\vect{\lambda},\tau,c^2,\sigma)
  \nonumber\\
  &\qquad=
  \left\{
    \prod_{i=1}^n
    p(y_i \mid \vect{\beta},\sigma,v_i)
  \right\}
  \left\{
    \prod_{i=1}^n
    p(v_i\mid \sigma)
  \right\}
  p(\vect{\beta}\mid \vect{\lambda},\tau,c^2)\,
  p(\vect{\lambda})\,p(\tau)\,p(c^2)\,p(\sigma),
\end{align}
with
\begin{equation}
  p(y_i \mid \vect{\beta},\sigma,v_i)
  =
  \frac{1}{\sqrt{2\pi\,\sigma B_q v_i}}
  \exp\!\left[
    -\frac{
      \big(
        y_i - \vect{x}_i^\top \vect{\beta} - A_q v_i
      \big)^2
    }{
      2\sigma B_q v_i
    }
  \right].
\end{equation}
Thus, up to an additive constant,
\begin{align}
  &\log p(\vect{y},\vect{v},\vect{\beta},\vect{\lambda},\tau,c^2,\sigma)
  \nonumber\\
  &\qquad=
  -\frac{1}{2}\sum_{i=1}^n
  \left[
    \log(\sigma B_q v_i)
    +
    \frac{
      \big(
        y_i - \vect{x}_i^\top \vect{\beta} - A_q v_i
      \big)^2
    }{
      \sigma B_q v_i
    }
  \right]
  \nonumber\\
  &\qquad\quad
  - \sum_{i=1}^n \left(\log \sigma + \frac{v_i}{\sigma}\right)
  - \frac{1}{2}\sum_{j=1}^p
  \left[
    \log V_j(\lambda_j,\tau,c^2)
    +
    \frac{\beta_j^2}{V_j(\lambda_j,\tau,c^2)}
  \right]
  \nonumber\\
  &\qquad\quad
  - \sum_{j=1}^p \log(1+\lambda_j^2)
  - \log\!\left(1+\frac{\tau^2}{\tau_0^2}\right)
  - \left(\frac{\nu}{2}+1\right)\log c^2
  - \frac{\nu s_0^2}{2c^2}
  + \log p(\sigma).
  \label{eq:logjoint-AL-RHS}
\end{align}

\subsection{Exact conditional kernel for the RHS scale block under AL}

Under the AL likelihood, conditional on $\vect{\beta}$, the RHS scale block again separates
from the likelihood:
\begin{equation}
  p(\vect{\lambda},\tau,c^2 \mid \vect{\beta},\vect{y},\vect{v},\sigma)
  \propto
  p(\vect{\beta}\mid \vect{\lambda},\tau,c^2)\,
  p(\vect{\lambda})\,p(\tau)\,p(c^2).
\end{equation}
Therefore,
\begin{equation}
  p(\vect{\lambda},\tau,c^2 \mid \vect{\beta},\vect{y},\vect{v},\sigma)
  \propto
  p(\vect{\lambda},\tau,c^2 \mid \vect{\beta},\vect{y},\vect{v},\vect{s},\sigma,\gamma)
\end{equation}
and the exact AL conditional kernel is identical to
\eqref{eq:rhs-kernel-exAL}--\eqref{eq:rhs-logkernel-exAL}. Equivalently, the transformed
log-target in \eqref{eq:rhs-logtarget-transformed} is unchanged.

In particular, the difference between the exAL and AL formulations lies in the likelihood
and in the auxiliary-latent structure $(\vect{v},\vect{s},\gamma)$, but \emph{not} in the
exact conditional algebra of the RHS latent scale block once $\vect{\beta}$ is held fixed.

\subsection{Full joint log-kernel under the AL likelihood}

Under the AL likelihood, the full joint log-kernel is obtained by dropping all additive
constants that do not depend on
\[
  (\vect{\beta},\vect{\lambda},\tau,c^2,\sigma,\vect{v}).
\]
Using the normal--exponential mixture representation of the AL likelihood, we have
\begin{align}
  &\log p(\vect{y},\vect{v},\vect{\beta},\vect{\lambda},\tau,c^2,\sigma)
  \nonumber\\
  &\qquad\propto
  -\frac{1}{2}\sum_{i=1}^n \log(\sigma B_q v_i)
  -\frac{1}{2}\sum_{i=1}^n
  \frac{
    \big(
      y_i-\vect{x}_i^\top\vect{\beta}-A_q v_i
    \big)^2
  }{
    \sigma B_q v_i
  }
  \nonumber\\
  &\qquad\quad
  -\sum_{i=1}^n \log \sigma
  -\sum_{i=1}^n \frac{v_i}{\sigma}
  \nonumber\\
  &\qquad\quad
  -\frac{1}{2}\sum_{j=1}^p \log V_j(\lambda_j,\tau,c^2)
  -\frac{1}{2}\sum_{j=1}^p \frac{\beta_j^2}{V_j(\lambda_j,\tau,c^2)}
  \nonumber\\
  &\qquad\quad
  -\sum_{j=1}^p \log(1+\lambda_j^2)
  -\log\!\left(1+\frac{\tau^2}{\tau_0^2}\right)
  -\left(\frac{\nu}{2}+1\right)\log c^2
  -\frac{\nu s_0^2}{2c^2}
  + \log p(\sigma).
  \label{eq:log-kernel-full-joint-AL}
\end{align}

If $\sigma \sim \IG(a_\sigma,b_\sigma)$ with kernel
\[
  p(\sigma) \propto \sigma^{-(a_\sigma+1)}
  \exp\!\left(-\frac{b_\sigma}{\sigma}\right), \qquad \sigma>0,
\]
then
\begin{align}
  &\log p(\vect{y},\vect{v},\vect{\beta},\vect{\lambda},\tau,c^2,\sigma)
  \nonumber\\
  &\qquad\propto
  -\frac{1}{2}\sum_{i=1}^n \log(\sigma B_q v_i)
  -\frac{1}{2}\sum_{i=1}^n
  \frac{
    \big(
      y_i-\vect{x}_i^\top\vect{\beta}-A_q v_i
    \big)^2
  }{
    \sigma B_q v_i
  }
  \nonumber\\
  &\qquad\quad
  -\sum_{i=1}^n \log \sigma
  -\sum_{i=1}^n \frac{v_i}{\sigma}
  \nonumber\\
  &\qquad\quad
  -\frac{1}{2}\sum_{j=1}^p \log V_j(\lambda_j,\tau,c^2)
  -\frac{1}{2}\sum_{j=1}^p \frac{\beta_j^2}{V_j(\lambda_j,\tau,c^2)}
  \nonumber\\
  &\qquad\quad
  -\sum_{j=1}^p \log(1+\lambda_j^2)
  -\log\!\left(1+\frac{\tau^2}{\tau_0^2}\right)
  -\left(\frac{\nu}{2}+1\right)\log c^2
  -\frac{\nu s_0^2}{2c^2}
  \nonumber\\
  &\qquad\quad
  -(a_\sigma+1)\log \sigma
  -\frac{b_\sigma}{\sigma}.
  \label{eq:log-kernel-full-joint-AL-IGsigma}
\end{align}

Combining the $\log \sigma$ terms yields the equivalent compact form
\begin{align}
  &\log p(\vect{y},\vect{v},\vect{\beta},\vect{\lambda},\tau,c^2,\sigma)
  \nonumber\\
  &\qquad\propto
  -\frac{1}{2}\sum_{i=1}^n \log(B_q v_i)
  -\left(\frac{3n}{2}+a_\sigma+1\right)\log \sigma
  \nonumber\\
  &\qquad\quad
  -\frac{1}{2}\sum_{i=1}^n
  \frac{
    \big(
      y_i-\vect{x}_i^\top\vect{\beta}-A_q v_i
    \big)^2
  }{
    \sigma B_q v_i
  }
  -\sum_{i=1}^n \frac{v_i}{\sigma}
  \nonumber\\
  &\qquad\quad
  -\frac{1}{2}\sum_{j=1}^p \log V_j(\lambda_j,\tau,c^2)
  -\frac{1}{2}\sum_{j=1}^p \frac{\beta_j^2}{V_j(\lambda_j,\tau,c^2)}
  \nonumber\\
  &\qquad\quad
  -\sum_{j=1}^p \log(1+\lambda_j^2)
  -\log\!\left(1+\frac{\tau^2}{\tau_0^2}\right)
  -\left(\frac{\nu}{2}+1\right)\log c^2
  -\frac{\nu s_0^2}{2c^2}.
  \label{eq:log-kernel-full-joint-AL-compact}
\end{align}

\end{document}
